% =====================================================================
%  Appendix A: Getting Started with R and Python
%  Source: IntroR.tex (Overleaf) — adapted
% =====================================================================

\chapter{Getting Started with R and Python}
\label{appx_Software}

\section{R and RStudio}

\textsf{R}\index{software!R} is an open-source implementation of the
object-oriented mathematical programming language \textsf{S}, developed by
statisticians worldwide. In contrast to commercial packages such as
\textsf{SPSS}\index{software!SPSS} or \textsf{Stata}\index{software!Stata},
\textsf{R} is a full mathematical programming language with capabilities
comparable to \textsf{MATLAB}\index{software!MATLAB}, not merely a program
that performs canned regression routines. It is free, has no licensing fees,
and is increasingly used by businesses, government agencies, and non-profit
organisations.

Run \textsf{R} via the integrated development environment (IDE)
\href{http://rstudio.org/}{RStudio}\index{software!RStudio} or via
\href{https://code.visualstudio.com/}{VS Code} with the R extension.

\subsection{Core Spatial Libraries for R}

\begin{lstlisting}[language=R, caption={Installing core spatial libraries}]
install.packages(c(
  "sf",          # vector spatial data
  "terra",       # raster spatial data
  "tidycensus",  # Census / ACS data API
  "tigris",      # TIGER/Line geographies
  "spdep",       # spatial weights and Moran's I
  "spatialreg",  # SAR, SEM, IV estimation
  "GWmodel",     # GWR
  "leontief",    # input-output modelling
  "fredr",       # FRED time series API
  "osmdata",     # OpenStreetMap
  "ggplot2",     # graphics
  "tidyverse"    # data manipulation
))
\end{lstlisting}

\subsection{A Minimal Working Example: Spatial Regression in R}

\begin{lstlisting}[language=R,
  caption={Spatial lag model with spdep and spatialreg}]
library(sf)
library(spdep)
library(spatialreg)

# Load data (any sf object with a response and covariates)
data(columbus)                     # bundled example dataset
col_sf <- st_as_sf(columbus)

# Construct spatial weight matrix (queen contiguity)
W_nb  <- poly2nb(col_sf)           # neighbour list
W_lw  <- nb2listw(W_nb)            # list-weight object

# OLS (ignores spatial dependence -- baseline)
ols <- lm(CRIME ~ INC + HOVAL, data = col_sf)

# Test for spatial autocorrelation in OLS residuals
moran.test(residuals(ols), W_lw)

# SAR model (spatial lag)
sar <- lagsarlm(CRIME ~ INC + HOVAL,
                data = col_sf, listw = W_lw)
summary(sar)

# SEM model (spatial error)
sem <- errorsarlm(CRIME ~ INC + HOVAL,
                  data = col_sf, listw = W_lw)
summary(sem)
\end{lstlisting}

\section{Python and the Spatial Stack}

Python provides an equally capable spatial analytics environment via the
PySAL ecosystem \citep{Rey2010}.

\begin{lstlisting}[language=Python,
  caption={Installing core spatial libraries (Python)}]
pip install geopandas libpysal esda spreg mgwr
\end{lstlisting}

\begin{lstlisting}[language=Python,
  caption={Spatial lag model with spreg (Python equivalent)}]
import geopandas as gpd
import libpysal
from spreg import GM_Lag

# Load data
data = gpd.read_file(libpysal.examples.get_path("columbus.shp"))

# Construct weights
W = libpysal.weights.Queen.from_dataframe(data)
W.transform = "r"  # row-standardise

# SAR model
y  = data[["CRIME"]].values
X  = data[["INC","HOVAL"]].values
sar = GM_Lag(y, X, w=W, name_y="CRIME",
             name_x=["INC","HOVAL"])
print(sar.summary)
\end{lstlisting}

\section{Quarto: The Reproducible Workflow}

Quarto unifies prose, code, and mathematics in a single document:

\begin{verbatim}
---
title: "Spatial Regression Analysis"
format: html
---

## Data and Weights

```{r}
#| echo: true
library(spdep)
data(columbus)
W_nb <- poly2nb(st_as_sf(columbus))
W_lw <- nb2listw(W_nb)
```

The Moran's $I$ test for the OLS residuals yields ...
\end{verbatim}

Render with \texttt{quarto render document.qmd}. Toggle between PDF and HTML
output by changing \texttt{format: pdf} or \texttt{format: html}.

\section{Getting Started with QGIS}

\textsf{QGIS}\index{software!QGIS} (\href{http://www.qgis.org}{qgis.org})
is a free, open-source desktop GIS that integrates with R and Python. For
most analytical tasks in this book, R and Python are preferable because
they produce reproducible, scriptable workflows. QGIS is most useful for:
(a) initial exploratory visual analysis of spatial data; (b) digitising
and editing geometries; (c) presentation-quality cartographic output.

\renewcommand\bibname{Further reading}
\begin{subappendices}
\bibliographystyle{econometrica}
\bibliography{Literature/OverLord}
\IfFileExists{Literature/AppA.bbl}{\input{Literature/AppA.bbl}}{\typeout{Need bibtex on AppA}}
\end{subappendices}
